\documentclass[12pt]{article}
\usepackage{amsmath}
\usepackage{geometry}
\geometry{a4paper, margin=1in}
\usepackage{graphicx}
\usepackage{hyperref}
\usepackage{booktabs}
\usepackage{listings}
\usepackage{xcolor}

\lstset{
  basicstyle=\ttfamily\footnotesize,
  keywordstyle=\color{blue},
  commentstyle=\color{gray},
  frame=single,
  breaklines=true,
  postbreak=\mbox{\textcolor{red}{$\hookrightarrow$}\space},
  captionpos=b
}


\title{ASHRAE 241 Risk Model Equations}
\author{}
\date{}

\begin{document}

\maketitle

\section*{Nomenclature}

\begin{tabular}{ll}
$\gamma$         & Surface deposition rate (h$^{-1}$) \\
$\lambda$        & Biological decay (h$^{-1}$) \\
$\mu$            & Face mask efficiency \\
$\overline{d}$   & Mean diameter of dehydrated respiratory particles (m) \\
$\phi$           & Total removal rate (h$^{-1}$) \\
$C_{drop}$       & Respiratory particles per unit volume of exhaled air (\#m$^{-3}$) \\
$D$              & Duration of exposure (h) \\
$DK$             & Dose constant \\
$E$              & Evaporation coefficient \\
$GVL_i$          & Genomic viral load (RNA copies m$^{-3}$) \\
$i$              & Individual person in a space \\
$I_0$            & Number of people in a space \\
$P(I)$           & Probability of a susceptible person becoming infected \\
$PBR$            & Pulmonary breathing rate (m$^3$h$^{-1}$) \\
$Q$              & Quanta of pathogen deposited in the respiratory tract of an exposed person \\
$QER_i$          & Quanta emission rate of a person (quanta h$^{-1}$ per person) \\
$R$              & Fractional reduction of viability after emission \\
$RTD$            & Respiratory tract deposition fraction \\
$T_{\text{ECA}_i}$ & Total equivalent clean airflow rate (m$^3$h$^{-1}$) \\
$V^{*}_{drop}$   & Expelled volume ratio of hydrated respiratory particles to air \\
$VER_i$          & Virion emission rate (virions h$^{-1}$ per person) \\
$VF$             & Viable fraction \\
$VOL$            & Space volume (m$^3$) \\
$ECA_i$          & Equivalent clean airflow rate per person (L s$^{-1}$ per person) \\
ASHRAE           & The American Society of Heating, Refrigerating, and Air-Conditioning Engineers \\
IRMM             & Infection risk management mode \\
RR               & Relative risk \\
WR               & Wells-Riley \\
\end{tabular}






\section*{Equations}

\begin{equation}
P(I) = 1 - \exp\bigl(-Q\bigr)
\tag{1}
\end{equation}

\begin{equation}
Q = \frac{PBR \cdot D \cdot (1 - \mu)^2}{\phi \cdot VOL} \sum_{i=1}^{I_0} QER_i
\tag{2}
\end{equation}

\begin{equation}
QER_i \equiv RTD \cdot VER_i / DK
\tag{3}
\end{equation}

\begin{equation}
VER_i = PBR \cdot V^{*}_{drop} \cdot GVL_i \cdot VF \cdot RD
\tag{4}
\end{equation}

\begin{equation}
V^{*}_{drop} = \frac{\pi}{6} \left( \overline{d}E \right)^3 C_{drop}
\tag{5}
\end{equation}

\begin{equation}
\phi = \gamma + \lambda + \frac{T_{\text{ECA}_i}}{VOL}
\tag{6}
\end{equation}

\begin{equation}
ECA_i = \frac{T_{\text{ECA}_i}}{I_0 \times 3.6}
\tag{7}
\end{equation}





\section*{Input Parameters}

\subsection*{Constants (Common to All Occupancy Categories)}

\begin{table}[h!]
\centering
\begin{tabular}{@{}lll@{}}
\toprule
\textbf{Parameter} & \textbf{Distribution} & \textbf{Value(s)} \\ \midrule
Viability loss fraction after exhalation, $R$ & Constant & 1.0 \\
Exposure duration, $D$ (h)                    & Constant & 1.0 \\
Community infection rate (\%)                & Constant & Healthcare: 3.0 \\
                                             &          & All other: 1.0 \\
Probability of infection, $P(I)$ (\%)        & Constant & 0.10 \\ \bottomrule
\end{tabular}
\caption{Input parameters common to all occupancy categories.}
\label{tab:input-constants}
\end{table}

\begin{lstlisting}[language=Python, caption={Python constants based on Table 1}]
R = 1.0   # Viability loss fraction after exhalation
D = 1.0   # Exposure duration (hours)
\end{lstlisting}



\subsection*{Variables (With Distributions)}

\begin{table}[h!]
\centering
\resizebox{\textwidth}{!}{%
\begin{tabular}{@{}lll@{}}
\toprule
\textbf{Variable} & \textbf{Distribution} & \textbf{Parameters / Range} \\ \midrule
Respiratory particle evaporation factor, $E$ & Beta & $\alpha = 5.0$, $\beta = 2.0$; min = 2.0, max = 5.0 \\
Genomic viral load, $GVL_i$ (log$_{10}$ RNAcopies ml$^{-1}$) & Normal & $\mu = 7.0$, $\sigma = 1.4$ \\
Viable fraction, $VF$ & Beta & $\alpha = 5.0$, $\beta = 2.0$; min = $10^{-4}$, max = $10^{-2}$ \\
Respiratory tract deposition fraction, $RTD$ & Uniform & min = 0.43, max = 0.65 \\
Dose constant, $DK$ & Uniform & min = 5.0, max = 15 \\
Biological decay rate, $\lambda$ (h$^{-1}$) & Log-normal & GM = 0.52, GSD = 1.9 \\
Surface deposition rate, $\gamma$ (h$^{-1}$) & Uniform & min = 0.42, max = 0.61 \\
\bottomrule
\end{tabular}%
}
\caption{Stochastic input variables and their distributions.}
\label{tab:input-variables}
\end{table}


\begin{flushleft}
\footnotesize
\textit{Note:} $GVL_i$ is converted into RNAcopies/m$^3$ by multiplying by $10^6$.\\
GM: geometric mean; GSD: geometric standard deviation; $\mu$: mean; $\sigma$: standard deviation; $\alpha$, $\beta$: shape parameters.
\end{flushleft}

\begin{lstlisting}[language=Python, caption={Python implementation of stochastic variable sampling}]
import numpy as np

# Respiratory particle evaporation factor (Beta distribution)
E = np.random.beta(a=5.0, b=2.0) * (5.0 - 2.0) + 2.0  # scaled Beta between 2.0 and 5.0

# Genomic viral load (Normal distribution)
GVL_log = np.random.normal(loc=7.0, scale=1.4)  # log10 RNAcopies/ml
GVL = 10**GVL_log * 1e6  # RNAcopies/m^3

# Viable fraction (Beta distribution)
VF = np.random.beta(a=5.0, b=2.0) * (1e-2 - 1e-4) + 1e-4

# Respiratory tract deposition fraction (Uniform distribution)
RTD = np.random.uniform(0.43, 0.65)

# Dose constant (Uniform distribution)
DK = np.random.uniform(5.0, 15.0)

# Biological decay rate (Log-normal distribution)
GM_lambda = 0.52
GSD_lambda = 1.9
lambda_ = np.random.lognormal(mean=np.log(GM_lambda), sigma=np.log(GSD_lambda))

# Surface deposition rate (Uniform distribution)
gamma = np.random.uniform(0.42, 0.61)
\end{lstlisting}


\subsection*{Variables by Occupancy Category}

\begin{lstlisting}[language=Python, caption={Occupancy-specific parameters without activity ratio}]
# Occupancy-specific parameters by first word of category
occupancy_params = {
    "Cell":     {"PBR_GM": 0.55, "PBR_GSD": 1.2, "Cdrop_GM": 14, "Cdrop_GSD": 1.1, "d_GM": 1.9, "d_GSD": 1.1},
    "Dayroom":  {"PBR_GM": 0.55, "PBR_GSD": 1.2, "Cdrop_GM": 19, "Cdrop_GSD": 1.1, "d_GM": 1.9, "d_GSD": 1.1},
    "Food":     {"PBR_GM": 0.55, "PBR_GSD": 1.2, "Cdrop_GM": 15, "Cdrop_GSD": 1.1, "d_GM": 1.9, "d_GSD": 1.1},
    "Gym":      {"PBR_GM": 0.62, "PBR_GSD": 1.3, "Cdrop_GM": 9.8, "Cdrop_GSD": 1.1, "d_GM": 1.8, "d_GSD": 1.1},
    "Office":   {"PBR_GM": 0.55, "PBR_GSD": 1.2, "Cdrop_GM": 14, "Cdrop_GSD": 1.1, "d_GM": 1.9, "d_GSD": 1.1},
    "Retail":   {"PBR_GM": 0.55, "PBR_GSD": 1.2, "Cdrop_GM": 11, "Cdrop_GSD": 1.1, "d_GM": 1.8, "d_GSD": 1.1},
    "Transportation": {"PBR_GM": 0.64, "PBR_GSD": 1.2, "Cdrop_GM": 15, "Cdrop_GSD": 1.1, "d_GM": 1.9, "d_GSD": 1.1},
    "Classroom": {"PBR_GM": 0.55, "PBR_GSD": 1.2, "Cdrop_GM": 17, "Cdrop_GSD": 1.1, "d_GM": 1.9, "d_GSD": 1.1},
    "Lecture": {"PBR_GM": 0.55, "PBR_GSD": 1.2, "Cdrop_GM": 9.8, "Cdrop_GSD": 1.1, "d_GM": 1.8, "d_GSD": 1.1},
    "Manufacturing": {"PBR_GM": 0.73, "PBR_GSD": 1.3, "Cdrop_GM": 15, "Cdrop_GSD": 1.1, "d_GM": 1.9, "d_GSD": 1.1},
    "Sorting": {"PBR_GM": 0.77, "PBR_GSD": 1.3, "Cdrop_GM": 19, "Cdrop_GSD": 1.1, "d_GM": 1.9, "d_GSD": 1.1},
    "Warehouse": {"PBR_GM": 0.77, "PBR_GSD": 1.3, "Cdrop_GM": 15, "Cdrop_GSD": 1.1, "d_GM": 1.9, "d_GSD": 1.1},
    "Exam": {"PBR_GM": 0.62, "PBR_GSD": 1.3, "Cdrop_GM": 32, "Cdrop_GSD": 1.1, "d_GM": 2.0, "d_GSD": 1.1},
    "Group": {"PBR_GM": 0.55, "PBR_GSD": 1.2, "Cdrop_GM": 12, "Cdrop_GSD": 1.1, "d_GM": 1.8, "d_GSD": 1.1},
    "Patient": {"PBR_GM": 0.62, "PBR_GSD": 1.3, "Cdrop_GM": 30, "Cdrop_GSD": 1.1, "d_GM": 2.0, "d_GSD": 1.1},
    "Resident": {"PBR_GM": 0.62, "PBR_GSD": 1.3, "Cdrop_GM": 21, "Cdrop_GSD": 1.1, "d_GM": 2.0, "d_GSD": 1.1},
    "Waiting": {"PBR_GM": 0.55, "PBR_GSD": 1.2, "Cdrop_GM": 11, "Cdrop_GSD": 1.1, "d_GM": 1.8, "d_GSD": 1.1},
    "Auditorium": {"PBR_GM": 0.55, "PBR_GSD": 1.2, "Cdrop_GM": 9.8, "Cdrop_GSD": 1.1, "d_GM": 1.8, "d_GSD": 1.1},
    "Place": {"PBR_GM": 0.50, "PBR_GSD": 1.2, "Cdrop_GM": 9.8, "Cdrop_GSD": 1.1, "d_GM": 1.8, "d_GSD": 1.1},
    "Museum": {"PBR_GM": 0.49, "PBR_GSD": 1.2, "Cdrop_GM": 9.8, "Cdrop_GSD": 1.1, "d_GM": 1.8, "d_GSD": 1.1},
    "Convention": {"PBR_GM": 0.49, "PBR_GSD": 1.2, "Cdrop_GM": 9.8, "Cdrop_GSD": 1.1, "d_GM": 1.8, "d_GSD": 1.1},
    "Spectator": {"PBR_GM": 0.55, "PBR_GSD": 1.2, "Cdrop_GM": 15, "Cdrop_GSD": 1.1, "d_GM": 1.9, "d_GSD": 1.1},
    "Lobbies": {"PBR_GM": 0.55, "PBR_GSD": 1.2, "Cdrop_GM": 9.8, "Cdrop_GSD": 1.1, "d_GM": 1.8, "d_GSD": 1.1},
    "Common": {"PBR_GM": 0.55, "PBR_GSD": 1.2, "Cdrop_GM": 9.8, "Cdrop_GSD": 1.1, "d_GM": 1.8, "d_GSD": 1.1},
    "Dwelling": {"PBR_GM": 0.051, "PBR_GSD": 4.3, "Cdrop_GM": 40, "Cdrop_GSD": 1.1, "d_GM": 2.2, "d_GSD": 1.1}
}

# Example usage
category = "Lecture"
params = occupancy_params[category]
print("Breathing rate GM for {}: {}".format(category, params["PBR_GM"]))
\end{lstlisting}





\subsection*{Volume of Occupancy Categories}

\begin{lstlisting}[language=Python, caption={Simplified occupancy categories and space volumes}]
# Simplified dictionary of space volumes by first word of occupancy category (m^3)
space_volumes = {
    "Cell": 320,
    "Dayroom": 600,
    "Food": 360,
    "Gym": 1900,
    "Office": 2700,
    "Retail": 6000,
    "Transportation": 10000,
    "Classroom": 320,
    "Lecture": 1800,
    "Manufacturing": 12000,
    "Sorting": 4800,
    "Warehouse": 1200,
    "Exam": 41,
    "Group": 270,
    "Patient": 81,
    "Resident": 81,
    "Waiting": 270,
    "Auditorium": 4600,
    "Place": 2400,
    "Museum": 10000,
    "Convention": 10000,
    "Spectator": 6300,
    "Lobbies": 4600,
    "Common": 4600,
    "Dwelling": 540
}

# Example: get the space volume by simplified category
category = "Lecture"
volume_m3 = space_volumes[category]
\end{lstlisting}





\end{document}
